% Options for packages loaded elsewhere
\PassOptionsToPackage{unicode}{hyperref}
\PassOptionsToPackage{hyphens}{url}
%
\documentclass[
  11pt,
]{article}
\usepackage{amsmath,amssymb}
\usepackage{lmodern}
\usepackage{iftex}
\ifPDFTeX
  \usepackage[T1]{fontenc}
  \usepackage[utf8]{inputenc}
  \usepackage{textcomp} % provide euro and other symbols
\else % if luatex or xetex
  \usepackage{unicode-math}
  \defaultfontfeatures{Scale=MatchLowercase}
  \defaultfontfeatures[\rmfamily]{Ligatures=TeX,Scale=1}
\fi
% Use upquote if available, for straight quotes in verbatim environments
\IfFileExists{upquote.sty}{\usepackage{upquote}}{}
\IfFileExists{microtype.sty}{% use microtype if available
  \usepackage[]{microtype}
  \UseMicrotypeSet[protrusion]{basicmath} % disable protrusion for tt fonts
}{}
\makeatletter
\@ifundefined{KOMAClassName}{% if non-KOMA class
  \IfFileExists{parskip.sty}{%
    \usepackage{parskip}
  }{% else
    \setlength{\parindent}{0pt}
    \setlength{\parskip}{6pt plus 2pt minus 1pt}}
}{% if KOMA class
  \KOMAoptions{parskip=half}}
\makeatother
\usepackage{xcolor}
\IfFileExists{xurl.sty}{\usepackage{xurl}}{} % add URL line breaks if available
\IfFileExists{bookmark.sty}{\usepackage{bookmark}}{\usepackage{hyperref}}
\hypersetup{
  pdftitle={Evaluating the Impact of Retrieval Strategies on the Accuracy of Retrieval-Augmented Generation Systems},
  hidelinks,
  pdfcreator={LaTeX via pandoc}}
\urlstyle{same} % disable monospaced font for URLs
\usepackage[margin=1in]{geometry}
\usepackage{longtable,booktabs,array}
\usepackage{calc} % for calculating minipage widths
% Correct order of tables after \paragraph or \subparagraph
\usepackage{etoolbox}
\makeatletter
\patchcmd\longtable{\par}{\if@noskipsec\mbox{}\fi\par}{}{}
\makeatother
% Allow footnotes in longtable head/foot
\IfFileExists{footnotehyper.sty}{\usepackage{footnotehyper}}{\usepackage{footnote}}
\makesavenoteenv{longtable}
\usepackage{graphicx}
\makeatletter
\def\maxwidth{\ifdim\Gin@nat@width>\linewidth\linewidth\else\Gin@nat@width\fi}
\def\maxheight{\ifdim\Gin@nat@height>\textheight\textheight\else\Gin@nat@height\fi}
\makeatother
% Scale images if necessary, so that they will not overflow the page
% margins by default, and it is still possible to overwrite the defaults
% using explicit options in \includegraphics[width, height, ...]{}
\setkeys{Gin}{width=\maxwidth,height=\maxheight,keepaspectratio}
% Set default figure placement to htbp
\makeatletter
\def\fps@figure{htbp}
\makeatother
\setlength{\emergencystretch}{3em} % prevent overfull lines
\providecommand{\tightlist}{%
  \setlength{\itemsep}{0pt}\setlength{\parskip}{0pt}}
\setcounter{secnumdepth}{-\maxdimen} % remove section numbering
\usepackage{float}
\usepackage{booktabs}
\hypersetup{colorlinks=true,linkcolor=black,citecolor=black,urlcolor=blue}
\ifLuaTeX
  \usepackage{selnolig}  % disable illegal ligatures
\fi
\usepackage[numbers]{natbib}
\bibliographystyle{unsrtnat}

\title{Evaluating the Impact of Retrieval Strategies on the Accuracy of
Retrieval-Augmented Generation Systems}
\usepackage{etoolbox}
\makeatletter
\providecommand{\subtitle}[1]{% add subtitle to \maketitle
  \apptocmd{\@title}{\par {\large #1 \par}}{}{}
}
\makeatother
\subtitle{CS391: Case Studies in Machine Learning, The University of
Texas at Austin}
\author{}
\date{November 25, 2026}

\begin{document}
\maketitle

\hypertarget{abstract}{%
\section{Abstract}\label{abstract}}

Retrieval-augmented generation can supply external evidence to a
language model, but better retrieval need not produce better answers.
This study compares BM25 keyword retrieval, pretrained semantic
retrieval, and a supervised logistic regression relevance reranker while
holding the generation model, questions, corpus, and scoring procedures
fixed. A reproducible Natural Questions prototype contains 500 training,
50 validation, and 150 test questions, with a shared corpus of 696
deduplicated source passages. The reranker combines BM25 score, semantic
similarity, question-passage word overlap, and passage length. Ten
conditions comprise an LLM-only baseline and each retrieval method at
K=3, 5, and 10, yielding 1,500 saved test responses from GPT-4.1-mini.
Every RAG condition exceeded the baseline's 22.00\% Exact Match and
0.4075 token F1. Reranker K=3 achieved the highest observed Exact Match,
48.67\%, and F1, 0.6656. Its paired F1 improvement over semantic K=3 was
+0.0321, with a pointwise 95\% bootstrap interval of {[}+0.0077,
+0.0618{]}. Corresponding differences at K=5 and K=10 had intervals
spanning zero. Larger K consistently increased input tokens but produced
mixed accuracy changes. Review of 24 deliberately selected questions
identified answer-format penalties, ambiguous references, and
source-present abstention. These findings support evaluating retrieval,
generation, and lexical scoring separately. The selected answer-bearing
corpus, limited test sample, and incomplete relevance annotations
constrain generalization.

\hypertarget{introduction}{%
\section{1. Introduction}\label{introduction}}

Language models are attractive question-answering interfaces because
they can generate short responses without requiring users to formulate
database queries. Their parameters can encode substantial relational
knowledge, although access varies across facts and relations
\citep{lama}. For an application that needs traceable evidence, a fluent
answer alone is insufficient: users also need to know which documents
were selected and whether those documents support the requested claim.
Retrieval-augmented generation introduces an external evidence source
rather than relying exclusively on parametric memory \citep{rag}.

This architecture creates a practical design question. Should a system
select passages through lexical matching, embedding similarity, or
supervised reranking? BM25 rewards overlapping terms and remains an
important retrieval baseline \citep{bm25}. Dense encoders can match
differently worded questions and passages \citep{dpr}. A reranker can
combine signals from both, but its additional computation and
supervision are useful only if improved ordering translates into better
downstream responses. These distinctions matter even when the language
model itself remains unchanged.

The project therefore studies retrieval as an experimentally controlled
component. It uses the same passage corpus, test questions, generation
snapshot, instructions, and answer metrics across conditions. A small
logistic model is trained to combine four observable question-passage
features, allowing a supervised machine-learning contribution without
training the language model or a neural cross-encoder. This setting
makes the comparison feasible for a course project while retaining
separate training, validation, and final evaluation stages.

A modest fixed corpus also makes passage-level disagreements
inspectable: the analysis can trace a changed response to saved source
IDs without reconstructing an evolving external search index.

Four questions organize the investigation. RQ1 asks how RAG compares
with an LLM-only baseline in answer accuracy. RQ2 asks how BM25,
semantic retrieval, and the trained reranker compare in source recovery
and downstream quality. RQ3 asks how K=3, 5, and 10 affect Exact Match,
token F1, token consumption, and recorded latency. RQ4 asks what
retrieval, generation, and evaluation failures help explain
disagreements. The final question requires examining actual evidence and
reviewer judgments, because a lexical mismatch cannot automatically
establish factual error.

The contribution is a reproducible empirical comparison rather than a
new general-purpose retrieval algorithm. It includes a fixed dataset and
corpus, a training-only logistic relevance model, 1,500 paired-condition
responses, an integrity audit, question-level bootstrap intervals, and a
systematic 24-question qualitative review. The central result is
conditional: semantic retrieval and reranking improved observed means
over BM25, but reranker gains depended on K, and additional context did
not consistently help.

\hypertarget{related-work}{%
\section{2. Related Work}\label{related-work}}

Open-domain question answering separates finding evidence from using it.
DrQA combined Wikipedia retrieval with a learned reader, showing why
document search and reading comprehension must be assessed together
\citep{drqa}. Original RAG combined parametric generation with a
retrieved nonparametric memory and jointly fine-tuned components
\citep{rag}. REALM incorporated retrieval into language-model
pretraining \citep{realm}, while latent retrieval research learned
retriever-reader behavior from answer supervision \citep{orqa}. These
architectures motivate evidence access, but they differ from the present
experiment: its generator is a fixed API model, and only a lightweight
relevance classifier is trained locally.

Sparse and dense retrieval encode different matching assumptions. The
probabilistic relevance framework explains BM25's term-based scoring
\citep{bm25}. Dense Passage Retrieval learns dual encoders for questions
and passages \citep{dpr}, whereas Sentence-BERT constructs reusable
sentence representations that support cosine comparisons \citep{sbert}.
The current semantic method uses a pretrained Sentence Transformer
without question-answer fine-tuning. This distinction prevents
attributing results to supervised dense-retriever training that was not
performed. Research comparing sparse, dense, and attentional
representations also identifies tradeoffs between lexical precision,
fixed-vector capacity, and richer interactions \citep{luan}.

Reranking revisits a candidate set rather than searching an unrestricted
collection from scratch. BERT passage reranking uses learned
query-passage interactions \citep{bert_rerank}; ColBERT retains
token-level late interaction while permitting document-side
precomputation \citep{colbert}. The present logistic model is
intentionally simpler: four scalar features replace contextual pair
encoding. Scikit-learn supplies consistent estimator interfaces for this
supervised pipeline \citep{sklearn_paper}. Its transparency supports
inspecting coefficient signs, although correlated features and class
weighting prevent causal feature-importance interpretations.

These alternatives also expose distinct training and inference
questions. A supervised dual encoder changes the representation space,
while a reranker changes the ordering of already available candidates.
Late interaction preserves finer matching information than a single
passage vector \citep{dpr, colbert, luan}. Comparing their published
headline scores would therefore mix architecture, supervision, corpus
and reader changes. The current experiment instead holds those external
conditions fixed where practical and asks whether a small learned
combination improves over its own candidate-generating baselines.
Because raw BM25 scores and cosine similarities use different scales,
feature standardization is part of the implemented classifier rather
than evidence that the underlying scores have identical meanings. This
narrower comparison contributes an inspectable local result without
claiming to reproduce the larger studies' training regimes.

Benchmark design determines what a retrieval comparison means. Natural
Questions contains real search queries paired with Wikipedia annotations
\citep{nq}. SQuAD established widely used span-oriented Exact Match and
token-overlap evaluation \citep{squad}, while its unanswerable extension
emphasized recognizing when evidence does not support an answer
\citep{squad2}. KILT evaluates provenance as well as downstream
predictions against a shared Wikipedia snapshot \citep{kilt}. BEIR
demonstrates that relative retrieval performance can change across
heterogeneous tasks and domains \citep{beir}. Consequently, results on a
selected answer-bearing corpus should not be presented as benchmark-wide
or cross-domain superiority.

Retrieved-context effects are similarly architecture dependent.
Fusion-in-Decoder showed that a trained generative reader could benefit
from additional passages \citep{fid}. In contrast, long-context studies
found sensitivity to evidence position \citep{lost_middle}, and
experiments with irrelevant arithmetic context demonstrated
distractibility \citep{distracted}. These results justify examining K
and evidence composition; they do not establish that any decrease here
was caused by distraction. The current experiment changes passage depth
but does not manipulate context positions or isolate individual
distractors.

Evaluation research further separates answer overlap from
evidence-supported factual precision. FActScore assesses atomic factual
support in long-form generation \citep{factscore}; that metric was not
implemented here. Instead, saved lexical scores are supplemented by
explicit human review. Statistical comparisons likewise require
attention to the evaluation unit. Bootstrap work in machine translation
resamples matched examples \citep{koehn}, and guidance for NLP
significance testing stresses matching methods to tasks and experimental
designs \citep{dror}. This study follows that principle by resampling
question IDs, preserving condition pairing, and reporting pointwise
uncertainty rather than treating 1,500 answers as independent questions.

\hypertarget{methodology}{%
\section{3. Methodology}\label{methodology}}

\hypertarget{dataset-and-preprocessing}{%
\subsection{3.1 Dataset and
preprocessing}\label{dataset-and-preprocessing}}

The implementation loads the \texttt{LLukas22/nq-simplified} repackaging
through Hugging Face Datasets \citep{nq_mirror, datasets}, pinned to
revision \texttt{d0b4aa2c6b142fda852d4f4f371c1134560b69fa}. This
derivative supplies question, context, and answer-text/character-offset
fields; it is not the original Natural Questions evaluation protocol.
Source training and test partitions were pooled in stable order. The
preparation manifest records 106,921 rows and 65 duplicate questions,
leaving 106,856 usable unique questions.

Questions, contexts, and answer annotations are validated, including
exact answer-offset correspondence. Deduplication uses Unicode NFKC
normalization, case folding, and word-token/whitespace normalization
while preserving original question text. Python's seeded random sampling
selects 700 questions from sorted normalized keys and partitions them
into 500 training, 50 validation, and 150 test questions. These custom
question splits remain disjoint. SHA-256 identifiers connect each
question to its retained source context. Exact context deduplication
yields 696 passages shared by all methods. Because the corpus is
assembled from selected sources across the custom splits, test source
texts are intentionally searchable; test question labels are withheld
from model fitting. This is a controlled answer-bearing prototype, not
an unbiased Wikipedia search collection.

\hypertarget{retrieval-approaches}{%
\subsection{3.2 Retrieval approaches}\label{retrieval-approaches}}

BM25 uses \texttt{rank\_bm25.BM25Okapi} \citep{rank_bm25}, lowercasing
and extracting Unicode word tokens without stemming or stopword removal.
Verified package defaults are k1=1.5, b=0.75, and epsilon=0.25; equal
scores retain passage-ID order. Semantic retrieval uses
\texttt{sentence-transformers/all-MiniLM-L6-v2}, a 384-dimensional
encoder whose default input limit is 256 wordpieces \citep{minilm_card}.
The saved model revision is
\texttt{1110a243fdf4706b3f48f1d95db1a4f5529b4d41}. Unit-normalized
float32 vectors are indexed in CPU FAISS \texttt{IndexFlatIP};
normalized inner products equal cosine similarity
\citep{faiss_paper, faiss_metrics}. Corpus vectors and passage IDs
retain identical row order. Embedding truncation affects representation,
while full original passages remain available to generation.

\hypertarget{trained-logistic-regression-reranker}{%
\subsection{3.3 Trained logistic regression
reranker}\label{trained-logistic-regression-reranker}}

Candidates are the deduplicated union of BM25 top-10 and semantic
top-10. Each pair receives BM25 score, semantic cosine similarity, the
fraction of unique question tokens occurring in the passage, and
full-passage word-token count. A pair is positive when its passage ID
equals the question's annotated source; other candidates are labeled
zero, meaning unannotated rather than necessarily irrelevant.

Training contains 8,794 pairs: 500 positives and 8,294 negatives. Twelve
otherwise-missing training positives were inserted; no positive
insertion occurred during validation or test retrieval. The pipeline
fits StandardScaler on training pairs and transforms later pairs with
those stored statistics \citep{scaler_docs}. LogisticRegression uses
balanced class weights, L2 regularization with the default lbfgs solver,
\texttt{max\_iter=1000}, and \texttt{random\_state=42}
\citep{logistic_docs}. C values 0.01, 0.1, 1, and 10 were compared on
validation MRR@10; C=0.01 was selected. Validation uses 853 natural
candidate pairs from 50 questions. The selected training-fitted pipeline
was saved without a train-plus-validation refit. Test candidates are
sorted by predicted label-1 probability; these class-weighted scores are
not calibrated factual-confidence estimates.

\hypertarget{rag-answer-generation}{%
\subsection{3.4 RAG answer generation}\label{rag-answer-generation}}

All conditions request \texttt{gpt-4.1-mini-2025-04-14}
\citep{openai_model}, temperature 0, a 64-output-token limit, and seed
42. Prompt version \texttt{nq-short-answer-v1} uses the same concise
short-answer system instruction. RAG user content contains the question
and ordered, numbered sources with passage IDs and complete text; it
asks for an evidence-based answer or ``Insufficient evidence.'' Baseline
questions omit documents and permit abstention. Reference answers and
source labels are not model inputs. A five-question validation pilot
preceded the frozen 150-question test. Usage, settings, retrieved IDs,
prompts, latency, retries, and responses were saved; successful calls
were cached. A single generation was obtained for each of ten conditions
per test question.

\hypertarget{evaluation-metrics}{%
\subsection{3.5 Evaluation metrics}\label{evaluation-metrics}}

Saved Exact Match applies lowercase, ASCII-punctuation removal,
English-article removal, and whitespace normalization, then maximizes
equality across accepted references. Token F1 maximizes
\texttt{2PR/(P+R)} across references, where precision and recall use
multiset token overlap, following the declared SQuAD-style policy
\citep{squad}. Empty or recognized abstained answers score zero. Hit@K
is one if the single annotated source occurs within rank K, otherwise
zero. MRR@10 averages reciprocal annotated-source rank, with zero when
absent from top-10. These measure annotated-source recovery, not
exhaustive relevance recall. Actual delivered-K source presence is
recorded separately. Efficiency uses provider-reported input/output
tokens and retrieval-plus-generation client milliseconds. Setup and
deliberate pacing are excluded; retries/backoff and network effects
remain. Monetary pricing was not configured.

\hypertarget{statistical-analysis-and-integrity}{%
\subsection{3.6 Statistical analysis and
integrity}\label{statistical-analysis-and-integrity}}

Every comparison pairs the same 150 question IDs. For score difference
\texttt{d(q)=score\_A(q)-score\_B(q)}, 10,000 seed-42 bootstrap samples
resample 150 question IDs with replacement; the same index matrix is
reused across comparisons. Percentile 2.5/97.5 limits describe mean
paired differences. EM effects are percentage points; F1 effects are
fractions. Intervals are pointwise, not multiplicity-adjusted, and do
not remove selection bias. Question exchangeability is assumed;
dependence from shared contexts is not separately modeled. The integrity
audit reported 47 passing checks, eight warnings, zero failures, and two
unverified checks. Frozen hashes, schemas, rankings, source mappings,
score agreement, and train-only provenance evidence were inspected; no
results were repaired for paper production.

\hypertarget{experimental-results}{%
\section{4. Experimental Results}\label{experimental-results}}

\hypertarget{overall-answer-accuracy}{%
\subsection{4.1 Overall answer accuracy}\label{overall-answer-accuracy}}

Table 1 reports all conditions; Figure 1 visualizes Exact Match. The
LLM-only baseline achieved 22.00\% EM and 0.4075 F1. Every RAG condition
had higher observed means. Across the nine RAG-versus-baseline
comparisons, both pointwise EM and F1 intervals excluded zero in the
positive direction. For reranker K=3, the EM increase was +26.67
percentage points, with interval {[}+18.67, +34.67{]}; the F1 increase
was +0.2581 {[}+0.1902, +0.3281{]}. This supports an improvement in
annotated-answer agreement on these questions, not a universal
factual-accuracy claim.

\begingroup\small
\renewcommand{\arraystretch}{1.15}
\begin{longtable}{@{}>{\raggedright\arraybackslash}p{(\linewidth - 12\tabcolsep) * \real{0.24}}
>{\raggedright\arraybackslash}p{(\linewidth - 12\tabcolsep) * \real{0.05}}
>{\raggedright\arraybackslash}p{(\linewidth - 12\tabcolsep) * \real{0.11}}
>{\raggedright\arraybackslash}p{(\linewidth - 12\tabcolsep) * \real{0.1}}
>{\raggedright\arraybackslash}p{(\linewidth - 12\tabcolsep) * \real{0.17}}
>{\raggedright\arraybackslash}p{(\linewidth - 12\tabcolsep) * \real{0.17}}
>{\raggedright\arraybackslash}p{(\linewidth - 12\tabcolsep) * \real{0.16}}@{}}
\caption{Answer scores and mean recorded usage/time; every
condition uses the same 150 questions.}\label{tab:results-1}\\
\toprule
Configuration & N & EM (\%) & F1 & Input tokens & Output tokens & Client
ms \\
\midrule
\endhead
LLM-only & 150 & 22.00 & 0.4075 & 160.8 & 7.25 & 581.2 \\
BM25 K=3 & 150 & 41.33 & 0.5697 & 1185.3 & 7.62 & 597.0 \\
BM25 K=5 & 150 & 39.33 & 0.5575 & 2193.4 & 8.35 & 1048.6 \\
BM25 K=10 & 150 & 42.67 & 0.5792 & 4528.7 & 8.67 & 662.0 \\
Semantic K=3 & 150 & 44.67 & 0.6336 & 1154.6 & 8.11 & 635.3 \\
Semantic K=5 & 150 & 47.33 & 0.6541 & 1651.6 & 8.37 & 636.8 \\
Semantic K=10 & 150 & 46.00 & 0.6492 & 3147.5 & 8.99 & 856.2 \\
Reranker K=3 & 150 & 48.67 & 0.6656 & 1208.5 & 8.30 & 662.4 \\
Reranker K=5 & 150 & 47.33 & 0.6578 & 1821.5 & 8.63 & 849.9 \\
Reranker K=10 & 150 & 46.00 & 0.6511 & 3352.5 & 8.75 & 719.1 \\
\bottomrule
\end{longtable}
\endgroup

\begin{figure}
\centering
\includegraphics[width=0.95\textwidth,height=\textheight]{../results/figures/figure_1_exact_match.pdf}
\caption{Mean Exact Match for the nine RAG conditions; the
dotted line represents the separate LLM-only baseline. Bars start at
zero. The highest observed result is highlighted descriptively, without
implying universal superiority.}
\end{figure}

\hypertarget{retrieval-quality}{%
\subsection{4.2 Retrieval quality}\label{retrieval-quality}}

Table 2 and Figure 2 compare shared top-10 rankings once per method.
Semantic retrieval and reranking recovered annotated sources more often
than BM25. Their MRR@10 values were 0.9302 and 0.9372, versus 0.7911 for
BM25. The reranker's natural union contained the source for 148 of 150
questions, setting a candidate-availability ceiling before reranking.
Source presence nevertheless did not ensure exact agreement: 639 of the
1,350 RAG records had a source in the delivered context but EM=0,
spanning 83 distinct questions. Source misses occurred in 106 records
across 28 questions. These overlapping configuration-level outcomes are
not independent question counts.

\begingroup\small
\renewcommand{\arraystretch}{1.15}
\begin{longtable}{@{}>{\raggedright\arraybackslash}p{(\linewidth - 10\tabcolsep) * \real{0.28}}
>{\raggedright\arraybackslash}p{(\linewidth - 10\tabcolsep) * \real{0.144}}
>{\raggedright\arraybackslash}p{(\linewidth - 10\tabcolsep) * \real{0.144}}
>{\raggedright\arraybackslash}p{(\linewidth - 10\tabcolsep) * \real{0.144}}
>{\raggedright\arraybackslash}p{(\linewidth - 10\tabcolsep) * \real{0.144}}
>{\raggedright\arraybackslash}p{(\linewidth - 10\tabcolsep) * \real{0.144}}@{}}
\caption{Annotated-source recovery from shared top-10 rankings;
all metrics are fractions.}\label{tab:results-2}\\
\toprule
Method & Hit@1 & Hit@3 & Hit@5 & Hit@10 & MRR@10 \\
\midrule
\endhead
BM25 & 0.7533 & 0.8200 & 0.8200 & 0.8600 & 0.7911 \\
Semantic & 0.9067 & 0.9400 & 0.9667 & 0.9800 & 0.9302 \\
Reranker & 0.9133 & 0.9600 & 0.9667 & 0.9800 & 0.9372 \\
\bottomrule
\end{longtable}
\endgroup

\begin{figure}
\centering
\includegraphics[width=0.95\textwidth,height=\textheight]{../results/figures/figure_3_retrieval_hit_at_k.pdf}
\caption{Annotated-source Hit@K and MRR@10 from saved full
rankings. Each method appears once because rankings are shared across
generation depths. These measures do not exhaust all relevant evidence.}
\end{figure}

\hypertarget{trained-reranker-comparison}{%
\subsection{4.3 Trained reranker
comparison}\label{trained-reranker-comparison}}

Reranker K=3 achieved the highest observed EM, 48.67\%, and F1, 0.6656.
Against BM25 K=3, its EM effect was +7.33 points {[}+2.67, +12.67{]},
with F1 effect +0.0960 {[}+0.0464, +0.1483{]}. Against semantic K=3, the
F1 effect was +0.0321 {[}+0.0077, +0.0618{]}. Table 3 and Figure 3 show
why that observation should not be generalized across K. At K=5 and
K=10, reranker and semantic EM means were equal; their small F1
differences had intervals spanning zero. K=3 question-level F1
wins/ties/losses were 30/110/10 against BM25 and 12/132/6 against
semantic, preserving unfavorable cases.

\begingroup\small
\renewcommand{\arraystretch}{1.15}
\begin{longtable}{@{}>{\raggedright\arraybackslash}p{(\linewidth - 6\tabcolsep) * \real{0.49}}
>{\raggedright\arraybackslash}p{(\linewidth - 6\tabcolsep) * \real{0.13}}
>{\raggedright\arraybackslash}p{(\linewidth - 6\tabcolsep) * \real{0.13}}
>{\raggedright\arraybackslash}p{(\linewidth - 6\tabcolsep) * \real{0.25}}@{}}
\caption{Selected paired effects; intervals use 10,000 question
resamples and are pointwise, without multiplicity correction.}\label{tab:results-3}\\
\toprule
Left minus right & EM change (pp) & F1 change & 95\% paired F1
interval \\
\midrule
\endhead
semantic K=3 minus bm25 K=3 & +3.33 & +0.0639 & {[}+0.0128,
+0.1180{]} \\
reranker K=3 minus bm25 K=3 & +7.33 & +0.0960 & {[}+0.0464,
+0.1483{]} \\
reranker K=3 minus semantic K=3 & +4.00 & +0.0321 & {[}+0.0077,
+0.0618{]} \\
reranker K=5 minus semantic K=5 & +0.00 & +0.0038 & {[}-0.0242,
+0.0339{]} \\
reranker K=10 minus semantic K=10 & +0.00 & +0.0019 & {[}-0.0226,
+0.0247{]} \\
reranker K=3 minus LLM-only & +26.67 & +0.2581 & {[}+0.1902,
+0.3281{]} \\
\bottomrule
\end{longtable}
\endgroup

\begin{figure}
\centering
\includegraphics[width=1\textwidth,height=\textheight]{../results/figures/figure_5_paired_f1_differences.pdf}
\caption{Selected paired mean F1 differences over the same 150
questions. Horizontal segments are saved percentile 95\% bootstrap
intervals from 10,000 resamples. Zero denotes no mean difference;
intervals are pointwise and not adjusted for multiple comparisons.}
\end{figure}

\hypertarget{influence-of-k}{%
\subsection{4.4 Influence of K}\label{influence-of-k}}

Figure 4 shows nonmonotonic F1. BM25 decreased from K=3 to K=5, then
increased at K=10; semantic retrieval peaked at K=5, and reranking
peaked at K=3. Reranker K=10 minus K=3 changed F1 by -0.0146, with
interval {[}-0.0411, +0.0107{]}. Semantic K=5 minus K=3 changed F1 by
+0.0205 {[}-0.0071, +0.0502{]}. Neither interval establishes a reliable
directional population effect. Across within-method contrasts, some
questions improved while others degraded, including cases where the
annotated source remained available.

\begin{figure}
\centering
\includegraphics[width=0.95\textwidth,height=\textheight]{../results/figures/figure_2_f1_vs_k.pdf}
\caption{Mean token F1 at K=3, 5, and 10, with the LLM-only
baseline shown separately. The axis spans 0--1. These are descriptive
means; paired-difference intervals are not misused as individual-mean
error bars.}
\end{figure}

\hypertarget{efficiency-and-latency}{%
\subsection{4.5 Efficiency and latency}\label{efficiency-and-latency}}

Figure 5 relates quality to input tokens. Reranker input usage rose from
1208.5 at K=3 to 3352.5 at K=10 while its mean F1 declined. All
within-method K increases consumed more input tokens. Table 1 also
reports short but nonidentical output-token means. Mean recorded total
client latency ranged from approximately 581 to 1,049 ms. BM25 K=5's
high mean and nonmonotonic times should be interpreted alongside two
recovered API errors and seven unusual wall-clock gaps. Outliers were
retained. Recorded retrieval means were approximately 4.5 ms for BM25,
27.5 ms for semantic, and 49.3 ms for reranking, the last including
candidate retrieval. No dollar estimate or strong intrinsic-speed
ranking follows from these noisy client measurements.

\begin{figure}
\centering
\includegraphics[width=0.95\textwidth,height=\textheight]{../results/figures/figure_4_f1_vs_input_tokens.pdf}
\caption{Mean F1 versus actual mean input tokens for all ten
configurations. Colors identify methods and labels distinguish K. Input
tokens are usage units, not estimated monetary costs. More context did
not consistently yield greater lexical answer quality.}
\end{figure}

\hypertarget{qualitative-failure-analysis}{%
\section{5. Qualitative Failure
Analysis}\label{qualitative-failure-analysis}}

The automated taxonomy separates source misses, source-present lexical
failures, method advantages/regressions, higher-K changes, and baseline
disagreements. Seed-42 selection combined three suggested starting
cases, category coverage, possible formatting mismatches, and randomized
fill to 24 unique questions. Every category had at least two eligible
selected questions. Categories overlap and the sample is purposeful, not
representative. All 24 primary comparisons were reviewed: 18 A responses
were labeled correct, five uncertain, and one partially correct.
Reviewer cause labels included nine lexical mismatches and six ambiguous
references. These judgments do not estimate factual accuracy over 150
questions or 1,500 responses.

Table 4 summarizes six primary comparisons. For the Friends ending,
semantic K=3 returned ``Ross Geller,'' while BM25 abstained. Semantic
and reranking retrieved the finale source first; BM25 omitted it. The
reviewer labeled A correct and supported. The observed advantage is
associated with source availability, although it does not isolate the
causal contribution of semantic matching or rule out useful alternative
evidence.

For Bull Run, BM25 returned the accepted ``Confederate forces,'' while
reranker K=3 produced a longer sentence containing that answer. Both
supplied the annotated source first. The reviewer labeled the reranked
answer correct and supported, with a lexical mismatch. Its EM=0 and
F1=0.3077 therefore do not demonstrate a factual regression. In
contrast, the broad-base population-pyramid response was reviewed as
partially correct: the reviewer noted supported core content plus
awkward additional wording. Its lower score is retained rather than
explained away as harmless formatting.

The Witcher comparison illustrates ambiguity without a source-ranking
improvement. Both K=5 prompts shared their first four passages,
including the source at rank 1. That source distinguishes a 1986 short
story from a 1990 collection. Semantic answered ``1986''; reranking
answered the accepted ``1990.'' The reviewer labeled A correct but
categorized the cause as ambiguous reference. Different secondary
context and interpretation remain plausible explanations, not
established mechanisms.

Source availability also failed to prevent abstention. For ``who sings
the song sorry that i let you down,'' BM25 retrieved the annotated
passage naming NF but abstained. The reviewer marked factual correctness
uncertain, evidence support supported, and cause abstention;
lyric-versus-title wording was a possible explanation. This is not
automatically a confirmed factually wrong assertion.

Finally, Hallmark channel answers demonstrate higher-K lexical
degradation with retained evidence. BM25 K=3 returned ``565 (HD),''
while K=5 added ``DirecTV channel''; the reference was ``565.'' F1
decreased from 0.6667 to 0.4000, although the reviewer labeled primary A
correct and supported. Additional passages accompanied the change, but
distraction was not experimentally isolated. Together, these cases
distinguish retrieved-source access, interpretation, response
formatting, and factual judgment.

\begingroup\small
\renewcommand{\arraystretch}{1.15}
\begin{longtable}{@{}>{\raggedright\arraybackslash}p{(\linewidth - 10\tabcolsep) * \real{0.2}}
>{\raggedright\arraybackslash}p{(\linewidth - 10\tabcolsep) * \real{0.16}}
>{\raggedright\arraybackslash}p{(\linewidth - 10\tabcolsep) * \real{0.08}}
>{\raggedright\arraybackslash}p{(\linewidth - 10\tabcolsep) * \real{0.16}}
>{\raggedright\arraybackslash}p{(\linewidth - 10\tabcolsep) * \real{0.08}}
>{\raggedright\arraybackslash}p{(\linewidth - 10\tabcolsep) * \real{0.32}}@{}}
\caption{Six reviewed primary comparisons; factual labels and
cause labels are the reviewer's judgments, not population estimates.}\label{tab:results-4}\\
\toprule
Question topic & A & F1 A & B & F1 B & Primary-A reviewer label \\
\midrule
\endhead
Friends ending & semantic K=3 & 1.0000 & bm25 K=3 & 0.0000 & correct;
other \\
Bull Run winner & reranker K=3 & 0.3077 & bm25 K=3 & 1.0000 & correct;
lexical mismatch \\
First Witcher book & reranker K=5 & 1.0000 & semantic K=5 & 0.0000 &
correct; ambiguous reference \\
Let You Down singer & bm25 K=3 & 0.0000 & semantic K=3 & 0.0000 &
uncertain; abstention \\
Hallmark channel & bm25 K=5 & 0.4000 & bm25 K=3 & 0.6667 & correct;
lexical mismatch \\
Broad-base pyramid & reranker K=5 & 0.1429 & bm25 K=5 & 0.2000 &
partially correct; lexical mismatch \\
\bottomrule
\end{longtable}
\endgroup

\hypertarget{discussion}{%
\section{6. Discussion}\label{discussion}}

RQ1 is answered positively within this prototype: every RAG
configuration improved mean lexical agreement over the no-retrieval
baseline, and the paired intervals support positive sample-level
differences. That result is narrower than proving that every retrieved
response was factually better. RAG prompts also explicitly constrained
generation to supplied evidence, whereas the baseline used parametric
knowledge; the comparison concerns these complete conditions rather than
an isolated mathematical retrieval effect.

For RQ2, dense retrieval's stronger annotated-source recovery is
consistent with matching meaning across wording differences
\citep{dpr, sbert}. The lightweight reranker improved observed K=3
ordering and downstream quality, but little separated it from semantic
retrieval at larger K. Its standardized semantic-similarity coefficient
was positive and larger than the other coefficients; word overlap was
negative conditional on the remaining features. Correlation, candidate
selection and class weighting make these associations unsuitable as
causal importance estimates. Without feature ablations, the separate
contribution of each feature is unmeasured.

For example, reranker K=3 included the annotated source in 144 prompts
but produced 73 exact matches. That difference is a lexical evaluation
gap, not evidence of 71 factual errors. Additional ranking supervision
cannot by itself correct reference wording. Candidate-pool coverage,
delivered evidence availability, and generated-answer agreement should
therefore remain separate diagnostics rather than interchangeable
measures of one system property.

RQ3 reveals a practical quality-efficiency tradeoff: input consumption
increased consistently, while accuracy did not. Trained generative
architectures can benefit from aggregating more passages \citep{fid},
but that does not guarantee the same behavior for this fixed prompted
model. Context position, unrelated material and ambiguity are candidate
explanations \citep{lost_middle, distracted}, not demonstrated causes
here. Selecting the best observed test K is descriptive reporting, not a
new tuning step or a validated deployment policy.

RQ4 identifies evaluation design as part of the explanation. Source
recovery, lexical span agreement and reviewer-assessed correctness can
disagree. The reviewed Muslim-population example is especially
cautionary: a table's share of the world's Muslim population was used as
the reference for a differently worded population-percentage question.
The reviewer retained uncertainty despite lexical agreement. The
pipeline preserved that annotation rather than modifying the benchmark
after observing test results. Lightweight RAG design should therefore
report candidate coverage, downstream scores, costs in observable units,
and evidence-backed error analysis together, while reserving
comprehensive factuality assessment for a separate evaluation
\citep{factscore}.

\hypertarget{limitations}{%
\section{7. Limitations}\label{limitations}}

The corpus contains 696 selected answer-bearing contexts, not
unrestricted Wikipedia. It makes relevant sources available by
construction and provides relatively few distractors. Only 150 test
questions were evaluated, and the custom split was created from pooled
source partitions. Earlier BM25-only test retrieval inspection means the
test set was not historically untouched, even though final generation
settings and classifier selection were frozen before the final
evaluation.

One annotated passage per question does not exhaust relevance. Negative
reranker labels can include useful evidence, and a source miss need not
mean total evidence failure. Exact Match and F1 measure lexical
agreement, not independently verified factual correctness. Some
references are incomplete, ambiguous or time sensitive. The single
reviewer examined a deliberately selected, overlapping sample; there is
no inter-rater reliability estimate or representative factual-accuracy
measurement.

Only one generation snapshot, embedding model and controlled corpus were
evaluated, with one response per condition rather than repeated trials.
Twenty-six provider backend fingerprints appeared despite fixed
requested settings. Recorded timing includes client/network and retry
effects; no pure inference benchmark or configured monetary-price
calculation was available. Pointwise bootstrap intervals were not
multiplicity-adjusted, and question resampling does not separately
account for shared-context dependence or generation variability.

Training provenance is supported by code, hashes, scaler statistics and
saved metadata, not independently signed historical fit logs. Local
response records do not independently attest provider billing or
internal execution. The audit therefore establishes consistency with
documented limitations, not absolute correctness. No neural reranker,
factual-verification model, answer-revision mechanism, context-position
intervention, feature ablation or cross-corpus generalization experiment
was conducted.

\hypertarget{conclusion}{%
\section{8. Conclusion}\label{conclusion}}

This study examined how retrieval strategy and passage count affect
question answering when the corpus, generator and evaluation questions
remain fixed. All nine RAG conditions exceeded the LLM-only baseline on
mean Exact Match and token F1. Semantic retrieval improved observed
means over BM25, while a training-only logistic relevance reranker
achieved the highest observed result at K=3: 48.67\% Exact Match and
0.6656 F1. Paired uncertainty supports its K=3 differences on this
sample, but not a consistent advantage over semantic retrieval at K=5 or
K=10.

Increasing K raised input-token consumption without monotonic
answer-quality gains. Source recovery and downstream correctness
therefore require separate evaluation. The reviewed cases further showed
why lexical failures can reflect wording, reference ambiguity,
abstention or partially correct elaboration rather than a uniform
factual-error mechanism. These observations support transparent
reporting of both favorable and unfavorable outcomes.

Future work could extend the corpus with independent distractors,
evaluate additional embeddings and cross-encoder rerankers, obtain more
comprehensive relevance judgments, and measure factuality through
semantic or multi-reviewer evaluation. Repeated generation trials would
help separate context effects from generation variability. These are
proposed extensions, not completed experiments. The present contribution
is a controlled, auditable comparison with explicit scope and
uncertainty, rather than a claim that one retrieval strategy universally
produces more accurate RAG systems.

  \bibliography{references}

\end{document}
